% !TEX root = ../Rulebook.tex

In this chapter the general rules will be explained that are valid for all tests. This chapter is seperated into the sections robot, arena environment, Service Areas and Objects.

\section{Robots}
The robots used for competition shall satisfy professional quality standards. The concrete definition of these standards is to be assessed by the TC, comprising aspects such as sturdy construction, general safety, and robust operation. It is not required that the robots are certified for industrial use to allow for a faster and easier implementation of research work. However, some requirements have to be met to ensure the safety of all participants of the competition. If there are even vague doubts about the eligibility of using particular designs, parts, or mechanisms, the team should consult the TC well in advance.

\subsection{Design and Constraints} \label{ssec:RobotDesignAndConstraints}
%The robots need to comply with certain size constraints. A robot, including all parts attached to it as used in the competition, must be able to move by itself into a configuration so that it fits into a box of side lengths 80 cm x 55 cm x 110 cm (length x width x height). If all the robot's parts, such as manipulator or anything able to protrude outside of the previously specified box, are fully extended, the system must still not exceed a box of side lengths 120 cm x 80 cm x 160 cm (length x width x height). The organizers may specify further constraints, such as weight limits. If a team would like to apply a robot with deviating robot dimensions, it should contact the TC. Exceptions for specific robots are possible in case of small differences.

\paragraph{Perks}
The following assumptions are made about the kind of robots used in the competition:

\begin{itemize}
	\item At least one of the robots used by a team is mobile and moves on wheels. No specific assumptions are made about the kinematic design, but the mobile robots should be able to move on basically flat, sufficiently firm surfaces. Aerial robots are not allowed in this competition. 
	\item Robots are equipped at least one manipulator and are able to handle the Objects defined in Table~\ref{tab:manipulation_objects}, Table~\ref{tab:new_objects1} and Table~\ref{tab:new_objects2}.
	\item The manipulator of the robot should be designed and mounted on the robot such that it can handle the Objects at different heights, usually between $0\si{\centi\meter}$ and $40\si{\centi\meter}$ above the floor.
	\item The robots use sensors to obtain information about their whereabouts in the environment and the task-relevant Objects. 
	Usual types of sensors used are 2D-LiDARs, (3D-) Cameras, RaDAR, Ultrasonic range finders and other. 
	Sensors must not harm humans due to their physical properties (e.g. no laser class 2 or above).

\end{itemize}

\paragraph{Size}
There are no constraints regarding the size and weight of the used robots, but they have to fit in the arena defined in section \ref{sec:ArenaDesign}. The minimum passage width is $80\si{\centi\meter}$. The used robots must be able to maneuver in that space.

\paragraph{Electrical}
The used batteries may not exceed a capacity of 500Wh. 300Wh of capacity is recommended. The maximum voltage allowed on the robotic system is 60 V DC.

\paragraph{Movement}
The maximum speed of the robots may not exceed 1.5 m/s. The robot should also be able to stop within a reasonable distance on concrete floor.

\paragraph{Actuator Types}
Electric, pneumatic, and hydraulic actuation mechanisms are permitted, provided that they are constructed and produced according to professional standards and meet safety constraints. Combustion engines and any kind of explosives are strictly forbidden. Robots may not pollute or harm their environment in any way, e.g. by loss of chemicals or oil, spilling liquids, or exhausting gases.

\paragraph{Connectivity} The TC may require that robots are equipped with a wireless communication device of some sort (e.g. 802.11n), in order to communicate task specifications to the robots. 

\subsection{Example Robots}

Figure~\ref{fig:example_robots} shows three examples how a robot suitable to the competition can look like. 

\begin{figure} [h!]
	\begin{center}
		\subfloat{\includegraphics[width=0.3\textwidth]{./images/robots/AutonOhm_Ohmn3.jpg}} \hfill
		\subfloat{\includegraphics[width=0.31\textwidth]{./images/robots/tyrolics2.jpg}} \hfill
		\subfloat{\includegraphics[width=0.3\textwidth]{./images/robots/robotto_bot_small.jpg}} 
	\end{center}
	\caption{Examples mobile robot platforms that can be used for \RCAW. Robots from the teams AutonOHM (Nuremberg-Germany), Tyrolics (Innsbruck-Austria) and robOTTO (Magdeburg-Germany) -- from left to right. }
	\label{fig:example_robots}
\end{figure}


\subsection{Behavior and Safety} \label{ssec:RobotBehaviorAndSafety}
For safety the robots have to meet the constraints in section \ref{ssec:RobotDesignAndConstraints}. In general, all robots shall be operated with maximum safety in mind. Any robot operation must be such that a robot neither harms humans nor damages the environment. 
The used batteries shall be handled with care and all team members must be educated in the correct usage, charging and storage of the batteries of the team. For lithium batteries appropriate storage bags must be used by the teams. The OC supplies a fire extinguisher for lithium batteries at the competition. If this is not sufficient for the used batteries of a team. The team is responsible for supplying an appropriate fire extinguisher by themselfs. The OC and TC control the observance of this rules.

All robots must have an emergency stop button. The emergency stop has to be a hard stop mechanism, that ensures that the energy transfer to all actuators is stopped immediately and the robot halts. The mechanism must be a red emergency stop button that is clearly visible, easily accessible and per wire attached to the robot. It has to be easy accessible from at least 3 sides of the robot. A wireless emergency stop button is optional but not sufficient.

The OC may request the proof of a robot's safety (e.g. the correct operation of an emergency stop) anytime during the competition and exclude teams that cannot satisfy safety requirements.

When participating in a competition, the team may operate the robot only in their own team area, in the arenas provided (possibly constrained by a schedule assigning periods of time for exclusive use of the arena by a team or a group of teams), and in any other areas designated by the organizers for robot operation. Any operation of robots outside of these areas, e.g. in public areas or emergency paths, require prior permission by the OC.

\textbf{Safety test procedure:}

Before the competition starts each robot has to perform a safety test procedure. This will be included in the competition schedule given by the OC.
\begin{itemize}
	\item Inspection of the robot platform (sharp edges and general construction)
	\item Description of the included safety systems by the team leader 
	\item Test of emergency stop while standing still
	\item Test of emergency stop during navigation 
	\item Test of emergency stop during manipulation
\end{itemize}


\clearpage

\section{Arena Environment}
\label{sec:ArenaDesign}
The competition is held in an arena resembling an example layout of industrial manufacturing facilities. In this Section different parts of the environment are explained.

\subsection{Floor}

The floor is made of some firm material. This includes among others floors made of concrete, screed, timber, plywood, chipboard, laminated boards, linoleum, PVC flooring, or carpet. Some examples are illustrated in Figure \ref{fig:example_floors}. Floors may neither be made of loose material of any kind (gravel, sand, or any material which may damage the functioning of the robot's wheels) nor may such material be used on top of the floor. Liquids of any kind are not allowed. The floor may have spots of unevenness of up to $1\si{\centi\meter}$ in any direction (clefts, rifts, ridges, etc.).


\begin{figure} [h!]
	\begin{center}
		\subfloat{\includegraphics[height = 2cm]{./images/general_rules/example_floor_1.jpg}} \hspace{0.1cm}
		\subfloat{\includegraphics[height = 2cm]{./images/general_rules/example_floor_2.jpg}} \hspace{0.1cm}
		\subfloat{\includegraphics[height = 2cm]{./images/general_rules/example_floor_3.jpg}} \hspace{0.1cm}
		\subfloat{\includegraphics[height = 2cm]{./images/general_rules/example_floor_4.jpg}} \hspace{0.1cm}
		\subfloat{\includegraphics[height = 2cm]{./images/general_rules/example_floor_5.jpg}}\\
		\subfloat{\includegraphics[height = 2cm]{./images/general_rules/example_floor_6.jpg}} \hspace{0.1cm}
		\subfloat{\includegraphics[height = 2cm]{./images/general_rules/example_floor_7.jpg}} \hspace{0.1cm}
		\subfloat{\includegraphics[height = 2cm]{./images/general_rules/example_floor_8.jpg}} \hspace{0.1cm}
		\subfloat{\includegraphics[height = 2cm]{./images/general_rules/example_floor_9.jpg}} \hspace{0.1cm}
		\subfloat{\includegraphics[height = 2cm]{./images/general_rules/example_floor_10.jpg}}
	\end{center}
	\caption{Examples of floors that can be used for \RCAW arenas.}
	\label{fig:example_floors}
\end{figure}



\subsection{Walls and Virtual Walls}
\label{subsec: Walls and virtual Walls}

The arena consists of outer and inner Walls used to build structures, create obstacles or function as protection barriers for teams and viewers. Walls may be either physical (plank) or virtual (red/white Tape). All walls (physical and virtual) have an infinitely height.
The arena is completely enclosed by Walls (both types possible), meaning robots are not allowed to exit the arena during a run. All types of Walls won't be changed during the competition. If the robot touches a Wall or Virtual Wall it results in a Major Collision.

The height of a physical Wall must be not less than $20\si{\centi\meter}$ and not more than $40\si{\centi\meter}$ (but will be seen as infinite high). Most Walls have a uniform main color (white), but may be enforced by metal (aluminum framework) and decorated with sponsor logos or ads.

Virtual Walls are made of red/white Tape and may never be crossed during a run. The arena can contain Walls and Virtual Walls inside.

\begin{figure} [h!]
\centering
\includegraphics[width= 0.8\textwidth ]{./images/general_rules/barrier_tapes_in_china15.jpg}
\caption{Example of a typical arena. The red/white Tape indicates a Virtual Wall. The yellow/black Tape indicates a Virtual Obstacle and the white ashlar-formed cartonages are Obstacles.}
\label{fig:walls_and_virt_walls}
\end{figure}


\subsection{Tapes}
\label{subsec:Tapes}

The red/white Tape (Tesa signal $5\si{\centi\meter}$ width) is considered as a Virtual Wall and has an infinite height. The red/white Tape is static and won't be changed during the whole competition. It will be used inside the arena and as an outer border of the arena.  Touching the red/white Tape is considered as a Major Collision.

\textbf{Markup Tape:}

Green electrical tape is considered as markup Tape. This tape can be used everywhere, where it is useful. It is intended as a marker for the referees and teams and not for the robot. Therefore the color may deviate, but the color is not red or yellow to guarantee a clear difference to the other tapes. The tape is used to mark the START and FINISH area. Furthermore the tape can be used to mark the position of tables (especially $0\si{\centi\meter}$) and Walls. The latter ones are useful for restoring the arena in case of a Major Collision.


\subsection{Tables}
\label{subsec: Tables}

Tables in the league typically measure $80\si{\centi\meter}$ wide (the side the robot approaches) and $50\si{\centi\meter}$ deep. These are the intended dimensions, but a margin of +/- $10\si{\centi\meter}$ is allowed to accommodate potential construction variations. In general terms, the table shall be big enough to contain at least one manipulation zone as described in section \ref{ssec:ManipulationZone}. In general two manipulation zones are included on one table (one zone on each side of the Table). The used table heights are $0\si{\centi\meter}$, $5\si{\centi\meter}$, $10\si{\centi\meter}$ and $15\si{\centi\meter}$. See further down in this section for more information about the $0\si{\centi\meter}$-Table. See figure \ref{fig:ws} for illustration. 
The tolerance for the heigth is $\pm 2 \si{\centi\meter}$. During a competition the Table size is only fixed in the margin of $\pm 2 \si{\centi\meter}$, because of arbitrary surfaces (see \ref{subsec:Arbitrary_Surfaces_and_Decoys} for an explanation of arbitrary surfaces). 
The table is closed between the bottom and the table top. That's why the Table can be used as a reference for navigation, when the height is sufficient for the robot. Note that the height may change slightly with arbitrary surfaces during a competition and thus the table can sometimes be visible to the laser scanners and sometimes not. 
 
\begin{figure} [h!]
	\begin{center}
		\includegraphics[height = 4cm]{./images/arena/tables_small.jpg}	
	\end{center}
	\caption{left: 0\si{\centi\meter}-Table with arbitrary surface (approx. 1\si{\centi\meter} thick); right: PPT-Table which is a standard 10\si{\centi\meter}-Table with cavities}
	\label{fig:ws}
\end{figure}

If a table has a height of $0\si{\centi\meter}$, a green electrical Tape (see \ref{subsec:Tapes}) will mark the area. A $0\si{\centi\meter}$-Table can be active or inactive. Active means, that the current test (see table \ref{fig:test_specifications_instance}) includes $0\si{\centi\meter}$-Tables.
If the table is active, the surrounded area may be covered by a white sheet of paper or arbitrary surface. The material is not fixed.
The OC is responsible to replace it in case of pollution or tears. To reduce the tear, it is removed when the $0\si{\centi\meter}$-Table is not active. If the floor is white or the table shall have an Arbitrary Surface, no cover needs to be installed.
$0\si{\centi\meter}$-Table may be crossed and does not count as a collision. If the laid out white Surface is moved, it is not a collision.
If the robot touches an Object while navigating, this will be handled as a Minor Collision. Examples for $0\si{\centi\meter}$-Table are shown in figure \ref{fig:0cmws}.



\begin{figure} [h!]
		\begin{center}
			\subfloat{\includegraphics[width=0.4\textwidth]{./images/arena/table_0_inactive_small.jpg}} \hspace{0.2cm}
			\subfloat{\includegraphics[width=0.4\textwidth]{./images/arena/table_0_arb_small.jpg}} 
		\end{center}
	\caption{left: $0\si{\centi\meter}$-Table which is considered as inactive or as an active table with arbitrary surface; right: $0\si{\centi\meter}$-Table which is considered as active with an arbitrary surface}
	\label{fig:0cmws}
\end{figure}

\todo[inline]{Limit to 10cm}


\subsection{Arena Layout}
\label{ssec:ArenaGeneral}

%\textbf{Layout:}
The arena is a static 2D environment consisting of Walls, Tables and Obstacles etc. with a size of atleast 10 m$^2$ and not more than 120 m$^2$. 
An example layout is shown in fig. \ref{fig:arena_example}.

\begin{figure} [h!]
	\centering
	\includegraphics[width= 1\textwidth ]{./images/general_rules/arena_example.jpg}
	\caption{An exemplary setup of a \RCAW environment.}
	\label{fig:arena_example}
\end{figure}

Layouts may include rooms and hallways to create more realistic scenarios.
Service Areas (see section \ref{sec:Service_Areas}) mark the locations for robots to perform tasks.
Each requested Service Area must be accessible via atleast one path of $80\si{\centi\meter}$ width.

\begin{figure} [h!]
	\centering
	\includegraphics[width= 0.8\textwidth ]{./images/general_rules/arena_map_annotated}
	\caption{Annotated 2D map of the environment in fig. \ref{fig:arena_example}}
	\label{fig:arena_map_annotated}
\end{figure}

Each competition has a new and unique layout designed by the actual TC members.
It should feature:
\begin{itemize}
	\item Area 10 m$^2$ - 120 m$^2$
	\item Minimum distance between arena elements at least $80\si{\centi\meter}$ 
	\item Widespread Service Areas entailing robot movements
	\item Multiple paths between Service Areas
	\item START and FINISH area
\end{itemize}

\clearpage
\textbf{START and FINISH area:}
\label{subsubsec: Start and Goal Area}

One or two parts of the arena are separated with marking Tape and considered as START and FINISH area. The START and FINISH area can be the same or two independent areas. The robot may leave the START area and enter the FINISH area only once. The START/FINISH area is leaved/entered when the robot completely cross the corresponding Marking Tape. In figure \ref{fig:tapeconfig} an exemplary Tape configuration is shown. As soon as the robot enters the FINISH area or re-enters the START area the run ends. The FINISH area counts as a Service Area and reaching the FINISH area gives additional points. 

\begin{figure} [h!]
	\begin{center}
		\includegraphics[width= 0.6\linewidth]{./images/arena/tabes.pdf}
	\end{center}
	\caption{Exemplary tape configuration: black lines are physical walls, red white dashed lines are virtual walls, yellow black dashed lines are barrier tape and green lines are marking tape}
	\label{fig:tapeconfig}
\end{figure}




\textbf{Referee Zone}

The arena is enclosed by a padded strip of 1m. This area is reserved for the referees to enable them to move freely while judging a robot's performance. See section \ref{sec: teams and roles}.


\clearpage

\section{Service Areas}
\label{sec:Service_Areas}
A Service Area indicates a location for a robot where tasks (e.g. picking or placing Objects) have to be performed.
\subsection{General} 
\label{subsec:Service_Areas_General}


Such a location is usually a Table with a flat white top (see fig. \ref{fig:arena_example}), commonly referred to as Workstation (only in a case of a normal Table), but can also be a Rotating Table, Shelf, Precise Placement station or any other type needed for a specific task.
In order to successfully reach a Service Area, robots must position themselves in front of the Service Area in a way
that allows detection of the Objects of interest and the robot has to stand still. 

\begin{figure} [h!]
	\centering
	\includegraphics[width= 0.8\textwidth ]{./images/general_rules/arena_service_area_free_space}
	\caption{Free Space in front of a Service Area}
	\label{fig:arena_service_area_free}
\end{figure}

The arena layout must define where the "front" of a table is.
Figure \ref{fig:arena_map_annotated} gives an example for the definition of the position of each Service Area, marking them as WSx (Workstation x). The orientation only indicates the direction of the Service Area. It does not specify the robot's heading, which may be chosen by teams according to their individual robot design.
Tables may includes Service Areas which can be used from both sides (see fig. \ref{fig:arena_example}), which  are defined as two seperate workstations (e.g. WS5 \& WS6).
This makes smaller physical layouts possible that still provide the ability to create complex navigation challenges due to the amount of locations to visit.


\subsection{Arbitrary Surfaces and Decoys}
\label{subsec:Arbitrary_Surfaces_and_Decoys}

In order to make the operating environment more realistic, the Service Areas may contain different kinds of Arbitrary Surfaces (Figure \ref{fig:ast_surface_example}) with industrial items as Decoys (Figure \ref{fig:ast_example}). The Decoys are limited in height to max. 10~cm. The width and depth of decoys are only limited by the size of the manipulation area. The Arbitrary Surfaces don't have to be fixed mounted at the Service Areas. Examples of Arbitrary Surfaces can be different wood pattern, grass, alufoil, plexiglas etc. The integration of Arbitrary Surfaces and Decoys into tests is according the table \ref{fig:test_specifications_instance}. If a team choose AprilTag objects in a test then no arbitrary surfaces will be used.


\begin{figure}[h!]
	\centering
	\subfloat[]{\includegraphics[width=.23333\textwidth]{./images/Arbitrary_Surface_rug}}
	\hspace{.05\textwidth}
	\subfloat[]{\includegraphics[width=.23333\textwidth]{./images/Arbitrary_Surface_wood}}
	\hspace{.05\textwidth}
	\subfloat[]{\includegraphics[width=.23333\textwidth]{./images/Arbitrary_Surface_grass}}\\
	\subfloat[]{\includegraphics[width=.23333\textwidth]{./images/Arbitrary_Surface_fake_wood_1}}
	\hspace{.05\textwidth}
	\subfloat[]{\includegraphics[width=.263333\textwidth]{./images/Arbitrary_Surface_fake_wood_2}}
	\caption{Examples of arbitrary surfaces used for Service Areas.}
	\label{fig:ast_surface_example}
\end{figure}

\begin{figure}[h!]
	\begin{center}
		\subfloat[]{\includegraphics[width=.375\textwidth]{./images/realisticWorkingDeskI.jpeg}}
		\hspace{.05\textwidth}
		\subfloat[]{\includegraphics[width=.375\textwidth]{./images/realisticWorkingDeskII.jpeg}}
	\end{center}
	\caption{Examplary configuration of the working desks including objects and decoys.}
	\label{fig:ast_example}
\end{figure}

Placement of Decoys and arbitrary surfaces are done by the referees. Decoys could be any kind of object (tools, hardware, small electrical devices) but also standard objects from the rule book defined in section~\ref{ssec:ManipulationObjects} can be used. Figure~\ref{fig:ast_example} shows some examples. 

\section{Objects} 
\label{ssec:ManipulationObjects}

The Objects in \RCAW include a wide range of Objects relevant in industrial applications of robotics. They eventually cover any raw material, (semi-)finished parts or products as well as tools and possibly operating materials required for manufacturing processes. Object types are selected by the TC and shall vary in complexity due to different shapes, colors and uses of objects. Currently there are three sets of Objects used in the competition: Basic, Advanced and April Tagged. The following sections provide details about those sets. The exact appearance of the object used in the competition can slightly vary due to availability, e.g. different coating colors for some standard parts.

\subsection{Basic Object Set}

The Basic Object Set includes standard profile rails, screws and nuts with various sizes and masses.
% (see table~\ref{tab:manipulation_objects}).

\newcommand{\imageView}[1]{\includegraphics[width=2cm, valign=c]{#1}}
\newcommand{\rowpadding}{0.4cm}
\setlength\extrarowheight{\rowpadding}


\begin{table}[h!]

\begin{tabular}{|c|c|c|c|m{8cm}|}
\hline
ID & Object & Symbolic Description & Mass & Details \\
\hline

\texttt{1} & \imageView{./images/F20_20_B.jpg} & \texttt{F20\_20\_B} & 49 g & Small aluminium profile \newline
 Coating/Colour: black anodized\newline
 Size: $20 \times 20 \times 100 mm$ \\ [\rowpadding]
\hline

\texttt{2} & \imageView{./images/F20_20_G.jpg} & \texttt{F20\_20\_G} & 49 g & Small aluminium profile \newline
 Coating/Colour: gray anodized\newline
 Size: $20 \times 20 \times 100 mm$ \\ [\rowpadding]
\hline

\texttt{3} & \imageView{./images/S40_40_B.jpg} & \texttt{S40\_40\_B} & 186 g & Large aluminium profile\newline
 Coating/Colour: black anodized\newline
 Size: $40 \times 40 \times 100 mm$ \\ [\rowpadding]
\hline

\texttt{4} & \imageView{./images/S40_40_G.jpg} & \texttt{S40\_40\_G} & 186 g & Large aluminium profile \newline
 Coating/Colour:  gray anodized\newline
 Size: $40 \times 40 \times 100 mm$ \\[\rowpadding]
\hline

\texttt{5} & \imageView{./images/M20_100.jpg} & \texttt{M20\_100} & 296 g & Screw\newline
 ISO4014, DIN 931, EU 24014 \newline
 Coating/Colour: blank, black burnished \newline
 Size: M$20\times 100$ \\ [\rowpadding]
\hline

\texttt{6} & \imageView{./images/M20.jpg} & \texttt{M20} & 56 g & Small nut\newline
 ISO4032, DIN934, EU 24032  \newline
 Coating/Color: blank, black burnished \newline
 Size: M20 \\ [\rowpadding]
\hline

\texttt{7} & \imageView{./images/M30.jpg} & \texttt{M30} & 217 g & Large nut\newline
 ISO4032, DIN934, EU 24032  \newline
 Coating/Color: blank, black burnished \newline
 Size: M30 \\ [\rowpadding]
 \hline
\end{tabular}
\caption{\RCAW Basic Object Set}
\label{tab:manipulation_objects}
\end{table}

\clearpage

\clearpage
\subsection{Advanced Object Set}


For 2023, the RoCKIn object set (see 2022 rulebook or previous ones) is replaced by a selection of better available parts for a drive train, as well as some tools (see Section~\ref{sec:new_objects}). The new object set was introduced in a technical challenge in 2022 and is called Advanced Object Set.

\begin{table}[h!]
	\begin{tabular}{|c|m{2cm}|c|c|m{8cm}|}
		\hline
		ID & Object & Symbolic Description & Mass & Details \\
		\hline
%%%%%%%%%%%%%
		\texttt{20} & \imageView{./images/newObjects/welle3d.jpg} \newline
		\imageView{./images/newObjects/welleSchematic.JPG}
		& \texttt{Axis2} & 180 g & Steel axis \newline
		Misumi: SFUB25-25-F28-P17-T15-S10-Q20 \newline
		Coating/Colour: blank, black burnished \newline
		Length: 68 mm \newline
		Diameter: 17mm, M20 \newline
		%see Figure~\ref{fig:welle2Schematic} \newline
		
		\href{https://de.misumi-ec.com/vona2/detail/110302635710/?CategorySpec=00000146753%3a%3ab%2cc}{Misumi} (visited Januar 2022)\\

			\hline
%%%%%%%%%%%%
		\texttt{21} & \imageView{./images/newObjects/bearing2.JPG} & \texttt{Bearing2} & 100 g & Bearing \newline 
		SKF YAR203-2F\newline
		Coating/Colour: gray \newline
		Useable with housing\newline
		\href{https://www.skf.com/sg/products/rolling-bearings/ball-bearings/insert-bearings/productid-YAR%20203-2F}{SKF}  (visited Januar 2022)\\
		\hline
%%%%%%%%%%%%				
		\texttt{22} & \imageView{./images/newObjects/housing.JPG} & \texttt{Housing} & 60 g & Housing \newline 
		SKF P40\newline
		Coating/Colour: gray \newline
		Useable with bearing\newline
		\textbf{Remark:} needs two hex socket screw M8x10 (ISO 4762, DIN 912) and two M8 nuts (ISO 4032, DIN 934) \newline
		\href{https://www.skf.com/sg/products/mounted-bearings/ball-bearing-units/pillow-block-ball-bearing-units/productid-P%2040}{SKF}  (visited Januar 2022)\\
		\hline
%%%%%%%%%%%%	
		\texttt{23} & \imageView{./images/newObjects/motor.JPG} & \texttt{Motor2} & 350g & Motor 755\newline
		Coating/Colour: gray \newline
		Size: $66.7 \times 42.0 \si{\milli\meter}$\newline
		Diameter: $d_{axis}=5\si{\milli\meter}$, $l_{axis}=10\si{\milli\meter}$ \newline
		\href{https://www.amazon.de/EsportsMJJ-12V-36V-3500-9000Rpm-Drehmoment-Hochleistungsmotor/dp/B075D85KVV}{Amazon}  (visited Januar 2022)\\
		\hline
%%%%%%%%%%%%					
		\texttt{24} & \imageView{./images/newObjects/FlangedResinCollar.JPG} & \texttt{Spacer} &  & Flanged Spacer\newline
		Misumi CLJHJ25-30-70  \newline
		Coating/Colour: white \newline
		Size: $70\si{\milli\meter}$\newline
		Diameter: $d_{inner}=25\si{\milli\meter}$, $d_{outer}=30\si{\milli\meter}$ \newline
		\href{https://us.misumi-ec.com/vona2/detail/110300236450/?curSearch=%7b%22field%22%3a%22%40search%22%2c%22seriesCode%22%3a%22110300236450%22%2c%22innerCode%22%3a%22%22%2c%22sort%22%3a1%2c%22specSortFlag%22%3a0%2c%22allSpecFlag%22%3a0%2c%22page%22%3a1%2c%22pageSize%22%3a%2260%22%2c%2200000042362%22%3a%22mig00000001500952%22%2c%2200000042368%22%3a%22b%22%2c%22jp000157843%22%3a%22mig00000000344081%22%2c%22jp000157846%22%3a%22mig00000001417174%22%2c%22jp000157851%22%3a%22mig00000000344088%22%2c%2200000334029%22%3a%2230%22%2c%2200000334032%22%3a%2270%22%2c%22typeCode%22%3a%22CLJHJ%22%2c%22fixedInfo%22%3a%22MDM0000085422111030023645020110476153310093415426696895%7c14%22%7d&Tab=preview}{Misumi}  (visited Januar 2022)\\
						\hline

\end{tabular}
\caption{Advanced Object Set - Part 1}
\label{tab:new_objects1}
\end{table}


\begin{table}[h!]
	\begin{tabular}{|c|m{2cm}|c|c|m{8cm}|}
		\hline
		ID & Object & Symbolic Description & Mass & Details \\
		\hline
		%%%%%%%%%%%%
		\texttt{25} & \imageView{./images/newObjects/weraScrewdriver.jpg} & \texttt{Screwdriver } & 19g & WERA 352 \newline
		Ball end screwdriver,hexagon socket screws\newline
		Coating/Colour: black/green \newline
		Size: $181\si{\milli\meter}$\newline
		Diameter: $d_{tip}=2.5\si{\milli\meter}$\newline
		Code: 05138070001\newline
		\href{https://www.amazon.de/Wera-05138070001-352-Sechskant-Kugelkopf-Schraubendreher-2-5/dp/B00154ZWFI?th=1}{Amazon}  (visited Januar 2022)\\
		\hline
		%%%%%%%%%%%%			
		\texttt{26} & \imageView{./images/newObjects/weraWrench.jpg} & \texttt{Wrench } & ca. 72g & WERA Jocker 6000, 8mm \newline
		Ratcheting combination wrenches\newline
		Coating/Colour: silver/grey/pink \newline
		Size: ca. $144\si{\milli\meter}$\newline
		Diameter: $d_{max}=20\si{\milli\meter}$\newline
		Code: 05073268001\newline
		\href{https://www.amazon.de/Wera-05073268001-Joker-Maul-Ringratschen-Schl%C3%BCssel/dp/B00BT0GBMG?th=1}{Amazon} (visited Januar 2022)\\
		\hline
		%%%%%%%%%%%%
		\texttt{27} & \imageView{./images/newObjects/hsscoDIN338metaldrill.jpg} & \texttt{Drill } & ca. 10g & Bosch Drill HSS-Co DIN338  \newline
		Drill\newline
		Coating/Colour: gold/ Cobalt alloy \newline
		Length: ca. $151\si{\milli\meter}$\newline
		Diameter: $d_{max}=13\si{\milli\meter}$\newline
		Code: 3165140382724 \newline
		\href{https://www.amazon.com/Bosch-2609255086-Metal-Drill-HSS-Co/dp/B0071OSFQY}{Amazon} (visited Januar 2022)\\
		\hline
		%%%%%%%%%%%
		\texttt{28} & \imageView{./images/newObjects/weraAllenKey.jpg} & \texttt{AllenKey } & ca. 10g & Wera Allen Key 8mm,\newline
		3950 PKL L-key, metric, stainless\newline
		Coating/Colour: silver \newline
		Length: ca. $195\si{\milli\meter} \times 37\si{\milli\meter}$\newline
		Diameter: $d_{max}=8\si{\milli\meter}$\newline
		Code: 05022708001 \newline
		\href{https://www.amazon.co.uk/Wera-WER022708-Hexagon-Keys-Multi-Colour/dp/B00A8QXTNG}{Amazon} (visited Januar 2022)\\
		\hline
		%%%%%%%%%%%
	
\end{tabular}
\caption{\RCAW Advanced Object Set - Part 2}
\label{tab:new_objects2}
\end{table}

\clearpage


\subsection{April Tagged Object Set}
\label{ssec: April Tagged Object Set}

For the season 2023 and onwards, 3D printed cubes marked with April tags might be used in the competition. 
This should allow teams to focus on other areas than object recognition by simplifying the detection of Objects. 
These cubes will be called April Tag Tagged Cubes (ATTC) in the following.

The April Tags used in \RCAW measure $40 \si{\milli\meter} \times 40 \si{\milli\meter}$ and have an encoding taken from the \textit{36h11} April Tag family, including a 1bit black and 1bit white border as shown in fig. \ref{fig:singleAprilTag}. 


\begin{figure}[h!]
	\centering
	\includegraphics[width= 0.4\textwidth ]{./images/singleAprilTAG42.png}
	\caption{Example of an April Tag with the ID 42. Encoding is from the 36h11-April Tag Family.}
	\label{fig:singleAprilTag}
\end{figure}

All the used April Tag ID numbers will fit to the ID numbers used by the AtWork Commander in the \href{https://github.com/robocup-at-work/atwork-commander/blob/master/atwork_commander_msgs/msg/Object.msg}{atwork\_commander\_msgs/msg/Object.msg}. 

Any software package to identify such april tags is allowed. The \href{http://wiki.ros.org/apriltag_ros}{April Tag ROS package} would be a great and easy-to-use example.

\subsubsection{Dimensions and Details of the ATTCs}

The standard ATTC measures $42\si{\milli\meter} \times 42\si{\milli\meter} \times 42 \si{\milli\meter}$ and has slightly rounded edges, similar to the shape of aluminum profile rails. 
A $1\si{\milli\meter}$ deep cavity measuring $40.5\si{\milli\meter} \times 40.5\si{\milli\meter}$ is embedded in the top of the cube to make centering the april tag printouts easier.
The exact design and size is specified in fig. \ref{fig:cubeObject} and can be downloaded from the official league repository: \href{https://github.com/robocup-at-work/rulebook/tree/alpha_2023/images}{3D STL file for an ATTC}

\begin{figure}[h!]
	\centering
	\subfloat[April Tag Tagged Cube (ATTC), including overall dimensions.]{\includegraphics[width=0.375\textwidth ]{./images/AprilTags/AprilTagCube2023_2.png}}
	\hspace{0.05\textwidth}
	\subfloat[Detailed dimensions of the ATTC cavaty to fix a certain April Tag.]{ \includegraphics[width=0.375\textwidth ]{./images/AprilTags/AprilTagCube2023_1.png} }
	\caption{3D construction of the ATTC as manipulation object.}%
	\label{fig:cubeObject}
\end{figure}

The cubes are 3D printed using PLA filament. Any color is allowed since the identification of cubes is solely done using the printout tag, which is to be glued into the top cavity.
Additionaly, printouts of a human-readable object ID and an image of the corresponding object from tables \ref{tab:manipulation_objects}, \ref{tab:new_objects1} and \ref{tab:new_objects2} must be glued to the side of the cubes to allow for easier onsite handling and identification of the cubes.

\begin{figure}[h!]
	\begin{center}
		\subfloat{\includegraphics[width=0.32\textwidth,rotate=-0]{./images/AprilTags/AT_obj_1.jpg}} \hfill
		\subfloat{\includegraphics[width=0.32\textwidth,rotate=-0]{./images/AprilTags/AT_obj_2.jpg}} \hfill
		\subfloat{\includegraphics[width=0.32\textwidth,rotate=-0]{./images/AprilTags/AT_obj_3.jpg}} \hfill
	\end{center}
	\caption{Examples of allowed ATTCs}
	\label{fig:real ATTCs}
\end{figure}



\subsubsection{Tagged Targets}

To allow for the more advanced tasks to be performed using the simplified april tag detection method, the target objects (precise placement cavities and containers) are also marked with a tag.

Containers are marked with a tag glued to the inside bottom of the container. The tag shall be placed as centered as possible.

For the precise placement tests, a special PP tile is used which allows a tag to be placed with a constant offset from the hole (see fig. \ref{fig:Tagged Targets}). The used ID matches the ID of the object that must be placed into that cavity.

\begin{figure}[h!]
	\begin{center}
		\subfloat{\includegraphics[width=0.32\textwidth,rotate=-0]{./images/AprilTags/AT_Container.jpg}} \hfill
		\subfloat{\includegraphics[width=0.32\textwidth,rotate=-0]{./images/AprilTags/CubeTile2023_2.png}} \hfill
		\subfloat{\includegraphics[width=0.32\textwidth,rotate=-0]{./images/AprilTags/CubeTile2023_1.png}} \hfill
	\end{center}
	\caption{Examples of tagged targets}
	\label{fig:Tagged Targets}
\end{figure}
